\section{The critical odd amplitude for profiles with quadratic factors}\label{sec:threshold}
Put $z=r^p$. The even and odd profiles below have respective factors
$(1-z)^2+\eta z$ and $(1-z)^2-qz$. For small $\eta>0$ the even
profile has a positive minimum of height asymptotic to $\eta$, and $q$
controls the amplitude of the perturbing odd monomial in $Q_q$.
We ask for which $q$ the first variation $J(H)$ takes a negative
value at some incoming height $-H$. Its strict signs transfer to finite
trajectories on each fixed finite list for sufficiently large dilation.

\begin{proposition}[Critical odd amplitude for the factors quadratic in $r^p$]\label{prop:threshold}
Fix even integers $d\ge10$ and $p\ge2$. For $0<\eta<1$ and $q>0$ let
\begin{equation}\label{eq:test-family}
 E_\eta=r^d\bigl((1-r^p)^2+\eta r^p\bigr),\qquad
 Q_q=r^{d-1}\bigl((1-r^p)^2-q r^p\bigr).
\end{equation}
Set $s_d=1-1/d$, $\alpha=3/2-s_d=1/2+1/d$ and $t_d=s_d/\alpha$.
Write $J(H)=A_\eta(H)-qB_\eta(H)$, with
\[
 (A_\eta(H),B_\eta(H))=
 \int_0^\infty \frac{\bigl(r^d(1-r^p)^2,\,r^{d+p}\bigr)}{(H+E_\eta(r))^2}\dd r,
 \qquad q_c(\eta)=\inf_{H>0}\frac{A_\eta(H)}{B_\eta(H)} .
\]
There is a depth $H>0$ with $J(H)<0$ exactly when $q>q_c(\eta)$.
As $\eta\downarrow0$,
\begin{equation}\label{eq:threshold-asym}
 q_c(\eta)\sim K_{d,p}\eta^\alpha,\quad
 K_{d,p}=\frac{2p c_d}{\pi}\frac{(1+t_d)^{3/2}}{t_d^{s_d}},
 \quad c_d=\int_0^\infty\frac{x^d\dd x}{(1+x^d)^2}.
\end{equation}
For sufficiently small $\eta$ a minimizing depth exists, and every
minimizing depth divided by $\eta$ tends to $t_d$.
If $\eta^\alpha=o(q)$ and $q=o(1)$, each of $V_o,V_i,J$ has an ordered
list of three strict signs $+,-,+$.
\end{proposition}

\begin{proof}
Abbreviate the two positive integrals by $A,B$. For each fixed
$\eta>0$, $B/A$ is bounded: off a neighborhood of $1$ the ratio of
the numerators is bounded, while on that neighborhood $E_\eta$ has a
positive lower bound. Its contribution to $B$ is $O_\eta((1+H)^{-2})$,
and integrating $A$ on a fixed subinterval of $(0,1)$ gives
$A\ge c_\eta(1+H)^{-2}$. Thus $0<q_c(\eta)<\infty$.
Since $B(H)>0$, the identity $J(H)=B(H)(A(H)/B(H)-q)$ gives
$J(H)>0$ at every depth for $q<q_c(\eta)$. At $q=q_c(\eta)$ it gives
$J(H)\ge0$, with equality exactly at a minimizing depth. For $q>q_c(\eta)$
some depth has $J(H)<0$.

Choose disjoint fixed small neighborhoods of $0$ and $1$. At zero,
$r=H^{1/d}x$ and dominated convergence give
$A_{\rm origin}\sim c_d H^{-s_d}$ as $H,\eta\to0$. A majorant after
scaling is $C x^d/(1+c x^d)^2$. Near $1$ we have
$E_\eta\asymp (r-1)^2+\eta$. There the numerator of $A$ is
$O((r-1)^2)$, giving $A_{\rm valley}=O((H+\eta)^{-1/2})$.
For $B$ put $r=1+\sqrt{H+\eta}\,x$. Uniformly for $H/\eta$ in compact
subsets of $(0,\infty)$, division of $H+E_\eta$ by $H+\eta$
gives $1+p^2 x^2+o(1)$ on bounded $x$. A majorant $C/(1+x^2)^2$
gives $B_{\rm valley}\sim \pi( H+\eta)^{-3/2}/(2p)$.
Off the valley $B$ is at most a constant times $A$, since
$r^p/(1-r^p)^2$ is bounded there. Away from both neighborhoods the
contribution to $A$ is uniformly bounded for small $H$ by positivity
and the integrable tail $O(r^{-d-2p})$. Since $1/2<s_d<1$,
\[
 \frac{A(t\eta)}{B(t\eta)\eta^\alpha}\ \longrightarrow \frac{2p c_d}{\pi}\,t^{-s_d}(1+t)^{3/2}
\]
locally uniformly for $t>0$.

To locate the infimum, for all small $H$ the origin gives
$A\ge c H^{-s_d}$ and the valley gives $B\le C A+C(H+\eta)^{-3/2}$.
For $H\ge h_0>0$ the same ratio $B/A$ is uniformly bounded: the
valley contribution to $B$ is $O(H^{-2})$, and an interval inside
$(0,1)$ bounded away from its endpoints gives $A\ge c(H+C)^{-2}$.
Consequently $\eta^\alpha B/A\to0$ along any sequence on which
$H/\eta\to0$ or $H/\eta\to\infty$. Continuity now confines minimizers
to a compact subinterval of $0<H/\eta<\infty$ for small $\eta$.
The limiting function has its unique minimum at $t_d$, proving
\eqref{eq:threshold-asym}. Thus for small $\eta$ the critical amplitude
is attained and $J$ vanishes at each minimizing depth when $q=q_c(\eta)$.
For $\eta^\alpha=o(q)$ one has $J(\eta)<0$.
At fixed $\eta,q$, as $H\downarrow0$ the origin gives
$J(H)\sim c_d H^{-s_d}$ since the extra power $r^p$ has vanishing
relative weight there. At $H=1$, dominated convergence as
$(\eta,q)\to(0,0)$ gives a strictly positive limit.

In $z=r^p$ the critical point equation for $E_\eta$ is
$d-(d+p)(2-\eta)z+(d+2p)z^2=0$. At $\eta=0$ its two simple roots
are $d/(d+2p)$ and $1$. The second persists as a minimum $c=1+O(\eta)$
of height $m=\eta+O(\eta^2)$, and the first gives a strictly higher
crest. The earlier inverse $E_\eta(a)=m$ satisfies $a\sim\eta^{1/d}$.
Near it $U(a)/f(a)\sim a^{d-1}/(2d-3)$ by the lowest powers.
Also $Q_q(c)=-q(1+O(\eta))+O(\eta^2)$. As $s_d>\alpha$,
$U(a)/f(a)+Q_q(c)<0$. Formula~\eqref{eq:Vo-Vi-band} makes both
positive-height tests negative just above $m$. Before that band the
lowest exponents give positive $U$ at sufficiently small height.
After it the highest exponents of $f,Q$ are $d+2p-2,d+2p-1$
with positive coefficients, so $U\to+\infty$, dominating the fixed
valley offset. This proves both remaining sign lists.
\end{proof}

\subsection{Ten-term realizations in degrees 22 and 26}
\begin{corollary}[Ten-term realizations at low degree]\label{cor:short-lists}
For each of the two supports listed below there is a nonempty open set
$V\subset(\mathbb R\setminus\{0\})^{10}$ such that every coefficient
vector in $V$ on that support gives at least sixteen simultaneous
hyperbolic limit cycles.
\[
 \{5,9,10,11,12,13,14,15,17,22\},\qquad
 \{5,9,10,13,14,17,18,19,21,26\}.
\]
\end{corollary}
\begin{proof}
Take the balanced seed on $\{5,9,10\}$ and set $d=10$, $q=\sqrt\eta$.
Since $\eta^{3/5}=o(q)$, Proposition~\ref{prop:threshold} supplies three
strict lists $(+,-,+)$ for both $p=2$ and $p=4$ when $\eta$ is
sufficiently small. For $p=2$ the valley adds $11,12,13,14$ and the
closure adds $15,17,22$. For $p=4$ the corresponding powers are
$13,14,17,18$ and $19,21,26$, giving
\begin{equation}\label{eq:portable-support}
 \{5,9,10,13,14,17,18,19,21,26\}.
\end{equation}
In both cases the minimum is simple and positive, and the leading
coefficients of the even and odd profiles are positive.
Lemma~\ref{lem:valley-rates} realizes the lists after inner dilation and
tail balance. Choose a common small physical shift. If $c>0$ is the
leading coefficient of the prepared primitive and $d_{\mathrm{top}}$ its
degree ($14$ or $18$), put $\gamma=c^{-1/(d_{\mathrm{top}}-1)}$. The
change $(x,y)=\gamma(X,Y)$ gives the monic primitive
$P_{A,\lambda}(\gamma X)/\gamma$.
For fixed $A$ the factor $\gamma$ is common to the balancing interval.
It divides every ordinate and the shift by the same number and retains
the adjustable odd coefficient. Proposition~\ref{prop:fold-closure}
now closes the lists.

The seed and the valley give four ordinary brackets. In decreasing incoming
height the lists $(+,-,+),(-,+),(+,-,+)$ precede the deep shifted minus
sign. Including the positive end signs gives eight incoming brackets.
The two fold brackets and two simple marks give four further cycles,
so Corollary~\ref{cor:sign-cost} yields $4+8+4=16$ simultaneous
hyperbolic limit cycles on each support. Hyperbolicity gives the asserted
open sets.
\end{proof}

\begin{remark}[Fixed-parameter verification of the low-degree profiles]\label{rem:fixed-low-profiles}
The proof of Corollary~\ref{cor:short-lists} uses the rational
degree-ten seed enclosures of Lemma~\ref{lem:compact-seed}\textup{(iii)}
and the small-$\eta$ threshold of Proposition~\ref{prop:threshold}.
At $\eta=2^{-20}$ and $q=2^{-10}$, $512$-bit Arb evaluation
\cite{Johansson2017} gives the same three local sign lists $(+,-,+)$
for $p=2$ and $p=4$. These fixed-parameter inequalities can also be used
as inputs to the dilation and matching proof.
Let $c$ be the minimum radius and $a$ its
upstream inverse at level $E(c)$. Bisection on the first rising branch
encloses $a$. Evaluating the polynomial $U$ and the moving-endpoint terms
in \eqref{eq:Vo-Vi-band} at the heights of $a/2,c+2^{-30},2$ gives
$(+,-,+)$ for both $V_o$ and $V_i$. Outward-rounded $512$-bit Arb bounds
for the negative-height integrals give:
\[
\begin{array}{c|ccc}
 p & J(\eta^2)/10^9 & J(\eta)/10^5 & J(1)\\ \hline
 2 & (6.99,7.02)&(-2.65,-2.60)&(0.0537,0.0538)\\
 4 & (6.9,7.1)&(-1.4,-1)&(0.044,0.045)
\end{array}
\]
For $p=2$ substitute $r=1/s$ on $[1,\infty)$. With
$B=(1-s^2)^2-qs^2$, $C=(1-s^2)^2+\eta s^2$, the entire integral is
\[
 \int_0^1 B\left(\frac{s^{10}}{(H+s^{10}C)^2}
                     +\frac{s^{12}}{(Hs^{14}+C)^2}\right)\dd s .
\]
On each of $2^{17}$ equal cells the two endpoints enclose both the
increasing powers of $s$ and the decreasing factor $(1-s^2)^2$.
For $p=4$, put $z=r^4$ and bound $E=r^{10}((1-z)^2+\eta z)$ and
$rQ=r^{10}((1-z)^2-qz)$ on $2^{16}$ equal cells on $[0,2]$.
Each lies on one side of $r=1$, so the endpoints again bound the square.
Sum the cell ranges times their widths. For $r\ge2$, $r^4-1\ge r^4/2$
gives $E\ge r^{18}/4$, and $|rQ|\le3r^{18}$ bounds the absolute
omitted integral by $48/(17\cdot2^{17})$.
In the stride-four profile the odd exponent $11$ is vacant, allowing
the support perturbation examined next.
\end{remark}

\begin{remark}[Sensitivity of the negative-height comparison]\label{rem:valley-insufficient}
The even minimum has height $\eta+O(\eta^2)$. For $q=\eta$ and sufficiently small
$\eta$, \eqref{eq:threshold-asym} gives $J(H)>0$ for every $H>0$.

To see the effect of changing the support, take $p=4$, $q=\sqrt\eta$
and consider the enlarged odd profile $Q_q+2q r^{d+1}$. Dividing by $r^{d-1}$
gives $(1-z)^2+q(2\sqrt z-z)$, $z=r^4$. It is strictly positive
for $0<z\le4$. For $z>4$ use $(z-1)^2/z\ge9/4$ and $q<1$.
Thus $J(H)>0$ at every depth for the enlarged profile. The added
coefficient tends to zero with $\eta$, and the even minimum, top degree
and two far matching directions are preserved.

For $d=10$ the added term in \eqref{eq:prepared-poly} is $2bqA^{-2}x^{11}$.
Compact convergence and the uniform tail bound following that equation
preserve the negative sign of the incoming seed test at the fixed
physical height $h_{\rm seed,1}$ after tail balance. Its scaled depth
$|h_{\rm seed,1}|/A^{10}$ tends to zero. Thus transfer of the
first-order signs at fixed depths gives no uniform sign for all actual
incoming heights.

In the undilated $d=10,p=4$ profile, the terminal even monomial changes
$E_\eta(1)$ to $\eta-(18/26)R^{-8}$. Positivity at this radius requires
$R^8\eta>18/26$. This necessary test holds for all sufficiently large
$R$ at fixed $\eta$ and fails if $R^8\eta\to0$ even as $R\to\infty$.
\end{remark}

\begin{remark}[Closure on other supports]\label{rem:portability}
Proposition~\ref{prop:fold-closure} starts with a balanced even-degree
primitive, a positive coefficient at the next odd power and an adjustable
lower odd coefficient. Once its bounded ordinary brackets and finite
shifted incoming signs are known, the same three exterior powers
close the list. The supplement applies the crest estimates to several
different inner supports. In Corollary~\ref{cor:short-lists} the
quadratic-factor profiles give three signs on each list. The repeated
gain of sixteen in Theorem~\ref{thm:main} comes from the cubic factors
\eqref{eq:isolated-pair} at stride $258$, which give seven, seven and
five signs. These finite lists give lower bounds. Finiteness and upper
bounds at a fixed $k\ge3$ remain open.
\end{remark}
